\section{Supplementary Results and Research Leads}
\label{app:supplementary-register-01}

This appendix keeps potentially relevant results from the research
archive attached to a mathematical home without enlarging the main
theorem spine.  This register is published separately from the six reader manuscripts.  Inclusion here is deliberately permissive.  Each entry
states its evidence boundary; entries labelled as open, contextual, or
evidence-only are not used as theorems elsewhere in the paper.  Earlier
formulations known to be false are replaced by their current correction.
The computational supplement and its \texttt{COMPUTATION.md} index give
the available program-level artifact map; an entry here does not turn a
shared-code check into an independent verification.
\begingroup\raggedright\emergencystretch=3em

\subsection{Additional results and constructions}

\phantomsection
\label{supp-note-01-001}
\paragraph{Three-dimensional Keller maps of every generic degree.}
\emph{Status.} Proof offered; independent review is not recorded.

For every integer d at least 3, there is a nonproper Keller map from complex affine 3-space to itself with generic degree d.

\emph{Credit.} Credit recorded to algal (construction, computation).

\phantomsection
\label{supp-note-01-002}
\paragraph{Degree-three function-field extension with S3 Galois closure.}
\emph{Status.} Proof offered; independent review is not recorded.

The function-field extension induced by the Alpöge map has degree 3 and an S3 Galois closure, as described by the cited explicit cubic model.

\emph{Credit.} Credit recorded to dorky (derivation).

\phantomsection
\label{supp-note-01-003}
\paragraph{The double-root resultant slice is affine three-space.}
\emph{Status.} Proof offered; independent review is not recorded.

In the binary-form multiplication model of the counterexample, the resultant-one double-root slice is polynomially isomorphic to affine three-space.

\emph{Credit.} Credit recorded to Alejandro Radisic (formalization); Andy Jiang (problem suggestion); ChatGPT (research assistance); Daniel Litt (proof); David Speyer (proof); Levent Alpöge (construction); Lillian Ryan Uhl (verification); Ravi Vakil (exposition); Terence Tao (proof, exposition); Will Sawin (proof).

\phantomsection
\label{supp-note-01-004}
\paragraph{For multiplication-incidence opens from P\textasciicircum{}a x P\textasciicircum{}b, the next (2, 3)...}
\emph{Status.} Proof offered; independent review is not recorded.

For multiplication-incidence opens from P\textasciicircum{}a x P\textasciicircum{}b, the next (2, 3) candidate fails to have the affine-space class/count expected of A\textasciicircum{}5; the original (1, 2) case is singled out within the stated incidence-coordinate class.

\phantomsection
\label{supp-note-01-005}
\paragraph{A rank-three affine-linear family whose full simple-root incidence...}
\emph{Status.} Proof offered; independent review is not recorded.

A rank-three affine-linear family whose full simple-root incidence is affine three-space lies in the tangent-but-not-osculating orbit and recovers the base map up to left-right equivalence.

\phantomsection
\label{supp-note-01-006}
\paragraph{Every generic-degree-three Keller map factors through a finite...}
\emph{Status.} Proof offered; independent review is not recorded.

Every generic-degree-three Keller map factors through a finite normal cubic cover whose trace-zero module is rank-two reflexive and locally free away from finitely many target points.

\phantomsection
\label{supp-note-01-007}
\paragraph{Flatness of the finite cubic normalization is equivalent to...}
\emph{Status.} Recorded result or structural lead; it is not promoted beyond the evidence described here and in the project archive.

Flatness of the finite cubic normalization is equivalent to vanishing of a finite defect module expressible through Ext data.

\emph{Credit.} Credit recorded to ChatGPT (AI-assisted derivation and drafting).

\phantomsection
\label{supp-note-01-008}
\paragraph{A minimal defect can only occur over an omitted singleton...}
\emph{Status.} Recorded result or structural lead; it is not promoted beyond the evidence described here and in the project archive.

A minimal defect can only occur over an omitted singleton fiber of scheme length at least four.

\phantomsection
\label{supp-note-01-009}
\paragraph{The defect problem admits a resolvent maximal-Cohen--Macaulay...}
\emph{Status.} Recorded result or structural lead; it is not promoted beyond the evidence described here and in the project archive.

The defect problem admits a resolvent maximal-Cohen--Macaulay reformulation.

\phantomsection
\label{supp-note-01-010}
\paragraph{At the first possible defect stratum the discriminant has...}
\emph{Status.} Recorded result or structural lead; it is not promoted beyond the evidence described here and in the project archive.

At the first possible defect stratum the discriminant has multiplicity at least six and the exceptional boundary has an elliptic-type constraint.

\phantomsection
\label{supp-note-01-011}
\paragraph{No algebraically homogeneous minimal defect can occur in a Keller...}
\emph{Status.} Proof offered; independent review is not recorded.

No algebraically homogeneous minimal defect can occur in a Keller normalization.

\phantomsection
\label{supp-note-01-012}
\paragraph{A generically degree-three Keller map admits a useful normalization...}
\emph{Status.} Recorded result or structural lead; it is not promoted beyond the evidence described here and in the project archive.

A generically degree-three Keller map admits a useful normalization picture: a finite normalization X over the target together with an open embedding of A3, expressible through a finite-flat binary-cubic model.

\phantomsection
\label{supp-note-01-013}
\paragraph{The choice A=C(C+1), B=-2-4C produces an exact cubic Keller family...}
\emph{Status.} Recorded result or structural lead; it is not promoted beyond the evidence described here and in the project archive.

The choice A=C(C+1), B=-2-4C produces an exact cubic Keller family with an unramified lost sheet, showing that omission need not coincide with ramification.

\phantomsection
\label{supp-note-01-014}
\paragraph{Within the A, B family, the nonproperness set can have arbitrarily...}
\emph{Status.} Recorded result or structural lead; it is not promoted beyond the evidence described here and in the project archive.

Within the A, B family, the nonproperness set can have arbitrarily many components, with discriminant and unramified-loss components controlled by the roots of explicit coefficient polynomials.

\phantomsection
\label{supp-note-01-015}
\paragraph{For a finite normalization of the cubic cover, the trace-zero...}
\emph{Status.} Recorded result or structural lead; it is not promoted beyond the evidence described here and in the project archive.

For a finite normalization of the cubic cover, the trace-zero module is reflexive and the nonflat locus is finite; flatness and the global cubic package require an additional Cohen-Macaulay hypothesis.

\phantomsection
\label{supp-note-01-016}
\paragraph{An affine-root incidence component of degree at least two cannot...}
\emph{Status.} Proof offered; independent review is not recorded.

An affine-root incidence component of degree at least two cannot be A2, and a generically squarefree affine-linear projective coefficient family has no connected full simple-root incidence isomorphic to A2.

\phantomsection
\label{supp-note-01-017}
\paragraph{The map is prime under polynomial composition.}
\emph{Status.} Proof offered; independent review is not recorded.

The map is prime under polynomial composition.

\phantomsection
\label{supp-note-01-018}
\paragraph{The off-diagonal collision space is a smooth factorial affine...}
\emph{Status.} Exact or computer-assisted certificate offered; see the computation index for its scope and independence boundary.

The off-diagonal collision space is a smooth factorial affine threefold with trivial Picard group and unit group modulo constants \(\mathbb Z\).

\phantomsection
\label{supp-note-01-019}
\paragraph{Escape rates are \(\varepsilon^{-1/2}\) near a smooth discriminant...}
\emph{Status.} Recorded result or structural lead; it is not promoted beyond the evidence described here and in the project archive.

Escape rates are \(\varepsilon^{-1/2}\) near a smooth discriminant point and \(\varepsilon^{-2/3}\) near the cusp, with more degenerate arcs allowing larger half-integral exponents.

\phantomsection
\label{supp-note-01-020}
\paragraph{Every nonzero constant combination of the inverse-Jacobian vector...}
\emph{Status.} Recorded result or structural lead; it is not promoted beyond the evidence described here and in the project archive.

Every nonzero constant combination of the inverse-Jacobian vector fields is incomplete.

\phantomsection
\label{supp-note-01-021}
\paragraph{Formal extension of the ruling can hold while Zariski-local...}
\emph{Status.} Recorded result or structural lead; it is not promoted beyond the evidence described here and in the project archive.

Formal extension of the ruling can hold while Zariski-local algebraization of every lift fails; explicit rational threefold and same-incidence nonalgebraizable lifts refute the proposed general algebraization lemma.

\phantomsection
\label{supp-note-01-022}
\paragraph{The resultant-one factor space is isomorphic to SL2 times A1\_gamma...}
\emph{Status.} Proof offered; independent review is not recorded.

The resultant-one factor space is isomorphic to SL2 times A1\_gamma, and its universal marked-root map to the simple-root binary-cubic locus is etale; ordering the roots gives an S3 torsor.

\phantomsection
\label{supp-note-01-023}
\paragraph{For every generic-degree-three Keller normalization, base change to...}
\emph{Status.} Proof offered; independent review is not recorded.

For every generic-degree-three Keller normalization, base change to the affine source splits the cubic algebra as S times S[eta]/(eta\textasciicircum{}2-D), implying flatness at every attained target point and a canonical marked-root form on the source.

\phantomsection
\label{supp-note-01-024}
\paragraph{If a generic-degree-three Keller map had smooth irreducible...}
\emph{Status.} Recorded result or structural lead; it is not promoted beyond the evidence described here and in the project archive.

If a generic-degree-three Keller map had smooth irreducible nonproperness hypersurface, then it would be surjective, the hypersurface would be isomorphic to A2, and its discriminant double cover would admit a connected etale C3-cover.

\phantomsection
\label{supp-note-01-025}
\paragraph{For any polynomial potential H(T, c), the marked-root and slope...}
\emph{Status.} Recorded result or structural lead; it is not promoted beyond the evidence described here and in the project archive.

For any polynomial potential H(T, c), the marked-root and slope equations b=r-H\_T(t, c), 2a=H(t, c)+tb have Jacobian r/2; composing with t=y+1/x, r=2/x, c=w(x, y)-x\textasciicircum{}3z has Jacobian -2 whenever the formulas extend polynomially, and generic degree deg\_T H.

\emph{Credit.} Credit recorded to ChatGPT (AI-assisted derivation and drafting).

\phantomsection
\label{supp-note-01-026}
\paragraph{On a sheet marked by a simple root t\_i of...}
\emph{Status.} Recorded result or structural lead; it is not promoted beyond the evidence described here and in the project archive.

On a sheet marked by a simple root t\_i of a cubic, x\_i=2/P'(t\_i)=2/[A(c)(t\_i-t\_j)(t\_i-t\_k)]; collision of the marked root therefore sends that affine source branch to infinity while the finite completion retains the ramification.

\emph{Credit.} Credit recorded to ChatGPT (AI-assisted derivation and drafting).

\phantomsection
\label{supp-note-01-027}
\paragraph{For every coprime normalized cubic-frame pair in the stated family...}
\emph{Status.} Proof offered; independent review is not recorded.

For every coprime normalized cubic-frame pair in the stated family, the inverse cubic is irreducible with nonsquare discriminant over C(a, b, c), so its generic Galois closure and geometric monodromy are S3.

\emph{Credit.} Credit recorded to ChatGPT (AI-assisted derivation and drafting).

\phantomsection
\label{supp-note-01-028}
\paragraph{Over A(c) nonzero, the normalized cubic discriminant complement is...}
\emph{Status.} Recorded result or structural lead; it is not promoted beyond the evidence described here and in the project archive.

Over A(c) nonzero, the normalized cubic discriminant complement is the product of the punctured c-line and the centered three-point configuration space; if A has s distinct roots, its fundamental group is F\_s times B\_3 and the permutation monodromy is the standard B\_3-to-S3 quotient.

\emph{Credit.} Credit recorded to ChatGPT (AI-assisted derivation and drafting).

\phantomsection
\label{supp-note-01-029}
\paragraph{Modulo triangular target shears, all polynomial root-slope frame...}
\emph{Status.} Proof offered; independent review is not recorded.

Modulo triangular target shears, all polynomial root-slope frame potentials admit the stated unique torus-weight normal form, with necessary and sufficient layer conditions at weights m<0, m=0, m=1, and m>=2.

\emph{Credit.} Credit recorded to ChatGPT (AI-assisted derivation and drafting).

\phantomsection
\label{supp-note-01-030}
\paragraph{A reduced frame potential of root degree n>=4 has leading...}
\emph{Status.} Proof offered; independent review is not recorded.

A reduced frame potential of root degree n>=4 has leading coefficient vanishing to order at least floor(n/2)+1 at the retained infinity point; the bound is sharp and yields component degrees (3n+4, 3n+3, 4) for even n and (3n+2, 3n+1, 4) for odd n.

\emph{Credit.} Credit recorded to ChatGPT (AI-assisted derivation and drafting).

\phantomsection
\label{supp-note-01-031}
\paragraph{For every n>=3 the displayed explicit frame potential defines a...}
\emph{Status.} Recorded result or structural lead; it is not promoted beyond the evidence described here and in the project archive.

For every n>=3 the displayed explicit frame potential defines a three-variable polynomial Keller map of generic degree n and determinant -2, attaining the sharp frame-degree bound.

\emph{Credit.} Credit recorded to ChatGPT (AI-assisted derivation and drafting).

\phantomsection
\label{supp-note-01-032}
\paragraph{For every admissible coprime cubic-frame pair, the homogenized...}
\emph{Status.} Proof offered; independent review is not recorded.

For every admissible coprime cubic-frame pair, the homogenized inverse cubic defines a smooth finite-flat Gorenstein triple cover of affine three-space with trivial rank-two Tschirnhausen bundle and distinguished homogeneous-root frame.

\emph{Credit.} Credit recorded to ChatGPT (AI-assisted derivation and drafting).

\phantomsection
\label{supp-note-01-033}
\paragraph{For the normalized cubic discriminant, with H=3A(c)t+B(c), the...}
\emph{Status.} Recorded result or structural lead; it is not promoted beyond the evidence described here and in the project archive.

For the normalized cubic discriminant, with H=3A(c)t+B(c), the conductor in C[c, t] is H\textasciicircum{}2 and the conductor ideal in the discriminant ring is generated by B\textasciicircum{}2-3Ab and 18Aa+Bb; transversely the singularity is the ordinary cusp C[H\textasciicircum{}2, H\textasciicircum{}3] inside C[H].

\emph{Credit.} Credit recorded to ChatGPT (AI-assisted derivation and drafting).

\phantomsection
\label{supp-note-01-034}
\paragraph{Every Keller map has an infinite-dimensional first-order left-right...}
\emph{Status.} Recorded result or structural lead; it is not promoted beyond the evidence described here and in the project archive.

Every Keller map has an infinite-dimensional first-order left-right automorphism space over the dual numbers, naturally identified with polynomial vector fields on the target; stabilization also adds inert diagonal affine automorphisms.

\emph{Credit.} Credit recorded to ChatGPT (AI-assisted derivation and drafting).

\phantomsection
\label{supp-note-01-035}
\paragraph{For positive boundary length, the hidden ordinary complex-point...}
\emph{Status.} Proof offered; independent review is not recorded.

For positive boundary length, the hidden ordinary complex-point automorphism kernel vanishes and Aut\_LR(G\_\{A, B\}) is the finite stabilizer of the decorated boundary scheme (Z\_A, B|Z\_A).

\emph{Credit.} Credit recorded to ChatGPT (AI-assisted derivation and drafting).

\phantomsection
\label{supp-note-01-036}
\paragraph{For every coprime admissible cubic-frame pair, the generic...}
\emph{Status.} Proof offered; independent review is not recorded.

For every coprime admissible cubic-frame pair, the generic function-field extension has no nontrivial deck transformation; ordinary source automorphisms over the identity target are therefore trivial.

\emph{Credit.} Credit recorded to ChatGPT (AI-assisted derivation and drafting).

\subsection{Open problems and unresolved assertions}

\phantomsection
\label{supp-note-01-037}
\paragraph{A failed general algebraization route.}
\emph{Status.} Open problem or unresolved assertion.

The proposed general algebraization lemma is false: explicit nonalgebraizable lifts provide counterexamples.  A narrower Keller-specific algebraization problem remains open.

\phantomsection
\label{supp-note-01-038}
\paragraph{Corrected divisorial taxonomy for a generic cubic normalization.}
\emph{Status.} Earlier question resolved at the divisorial level; independent
review is not recorded.

The complete list at a geometric generic target divisor is
\(U_0,U_1,U_2,B\), where \(U_i\) has trivial inertia and exactly \(i\)
deleted sheets, and \(B\) has transposition inertia with the ramified point
deleted and the unramified point retained.  The earlier \(U_1/U_2/B\) list
omitted the ordinary all-retained case \(U_0\).  This classification is
divisorial and does not control isolated nonflat points.

\phantomsection
\label{supp-note-01-039}
\paragraph{Can the common-root integral-closure theorem and left-right...}
\emph{Status.} Open problem or unresolved assertion.

Can the common-root integral-closure theorem and left-right equivalence result be independently proved or formalized?

\phantomsection
\label{supp-note-01-040}
\paragraph{The nonsquarefree common-root quotient.}
\emph{Status.} Open problem or unresolved assertion.

The later proof draft answers the earlier question: the full Artinian class B modulo A is retained at repeated roots.  Independent verification remains desirable.

\subsection{Context, corrections, and evidence boundaries}

\phantomsection
\label{supp-note-01-041}
\paragraph{Normal cubic covers with S3 monodromy need not be flat...}
\emph{Status.} Corrected statement; the earlier stronger formulation is not adopted.

Normal cubic covers with S3 monodromy need not be flat: an explicit minimal nonflat normal cubic algebra supplies a countermodel outside the Keller setting.

\phantomsection
\label{supp-note-01-042}
\paragraph{The remaining minimal defect can be organized by a radial...}
\emph{Status.} Recorded result or structural lead; it is not promoted beyond the evidence described here and in the project archive.

The remaining minimal defect can be organized by a radial fibration and an E6-type boundary/discrepancy picture.

\endgroup
